\chapter{The Unified Movement Budget: A Universal Conservation Law}

\author{James Tyndall}

%----------------------------------------------------------------------
\section{Introduction: Beyond Static Geometry to Dynamic Conservation}
%----------------------------------------------------------------------

The axioms of SDT describe a universe of tangible, geometric entities.  While this static picture is powerful, it is incomplete.  To describe a universe of motion, interaction, and change, we must introduce a dynamic principle.

That principle is the \textbf{Unified Movement Budget}.


%----------------------------------------------------------------------
\section{Deriving the Electron's Intrinsic k-Factor}
%----------------------------------------------------------------------

\subsection{The Vortex Energy Model}

A particle's rest energy ($E = mc^2$) is equivalent to the total rotational energy of its displacement vortex.

\subsection{The Universal Particle Geometry}

All stable fundamental particles (baryons and leptons) share a self-similar internal geometry.  The ratio of effective energetic volume to geometric volume is a universal constant, derived from the proton's structure:
\begin{equation}\label{eq:ratio}
  \text{Ratio}_{\text{particle}} = \frac{V_{\text{effective}}}{V_{\text{geometric}}} \approx 0.04
\end{equation}

The k-factor follows from the Vortex Energy Model:
\begin{equation}\label{eq:k-derive}
  k = \sqrt{7.5 \times \text{Ratio}_{\text{particle}}}
\end{equation}

\subsection{The Electron k-Factor}

\begin{equation}
  k_e = \sqrt{7.5 \times 0.04} = \sqrt{0.3} \approx 0.548
\end{equation}

This is not an empirical fit; it is a direct consequence of the universal geometry shared with the proton.


%----------------------------------------------------------------------
\section{The Conservation Law}
%----------------------------------------------------------------------

\subsection{The Movement Budget}

\begin{equation}\label{eq:budget}
  v_{\text{budget}} = k_e \cdot c \approx 0.548\,c
\end{equation}

This budget is the total, conserved, characteristic velocity of the electron's standing-wave vortex.

\subsection{Budget Allocation}

The budget is allocated among all orthogonal modes of motion:
\begin{equation}\label{eq:conservation}
  \boxed{v_{\text{orbital}}^2 + v_{\text{linear}}^2 + v_{\text{magnetic}}^2 = v_{\text{budget}}^2}
\end{equation}

\begin{itemize}
  \item $v_{\text{orbital}}$: internal ``orbital'' speed of the vortex wave pattern.  Defines the electron's quantum energy level and atomic clock tick rate.
  \item $v_{\text{linear}}$: external, linear velocity of the particle through the spation medium.
  \item $v_{\text{magnetic}}$: precessional velocity of the vortex spin axis aligning with an external magnetic field.
\end{itemize}

Every interaction is a \textbf{reallocation} of this finite, conserved budget.


%----------------------------------------------------------------------
\section{Deriving Relativistic Effects: Time Dilation}
%----------------------------------------------------------------------

An atomic clock's frequency ($f$) is determined by electron transition energies, which are a function of $v_{\text{orbital}}^2$.

\textbf{Condition:} Atom moving with $v_{\text{linear}} > 0$, no magnetic field ($v_{\text{magnetic}} = 0$).

\textbf{Mechanism:} To accommodate linear motion, the electron must divert budget from internal motion.  $v_{\text{orbital}}$ decreases.

\textbf{Derivation:}
\begin{align}
  v_{\text{orbital,moving}}^2 &= v_{\text{budget}}^2 - v_{\text{linear}}^2 \\
  \frac{f_{\text{moving}}}{f_{\text{rest}}} &= \frac{v_{\text{orbital,moving}}^2}{v_{\text{orbital,rest}}^2}
    = \frac{v_{\text{budget}}^2 - v_{\text{linear}}^2}{v_{\text{budget}}^2}
\end{align}

\begin{equation}\label{eq:time-dilation}
  \boxed{f_{\text{moving}} = f_{\text{rest}} \left(1 - \frac{v_{\text{linear}}^2}{k_e^2\,c^2}\right)}
\end{equation}

This is the \textbf{SDT Time Dilation formula}.  It provides a mechanical cause for the observed effect and makes a novel prediction: the magnitude of time dilation is scaled by $k_e^2$, suggesting a slight, measurable deviation from the standard relativistic formula.


%----------------------------------------------------------------------
\section{Deriving Quantum Effects: The Zeeman Effect}
%----------------------------------------------------------------------

\textbf{Condition:} Atom at rest ($v_{\text{linear}} = 0$), in external magnetic field $B$, forcing precession with $v_{\text{magnetic}} > 0$.

\textbf{Mechanism:} Precessional motion consumes budget, forcing $v_{\text{orbital}}$ to decrease.  This reduction in orbital kinetic energy IS the Zeeman energy shift.

\textbf{Derivation:}
\begin{align}
  v_{\text{orbital,field}}^2 &= v_{\text{budget}}^2 - v_{\text{magnetic}}^2 \\
  \Delta E_{\text{Zeeman}} &= \tfrac{1}{2}m_e\!\left[v_{\text{orbital}}(B{=}0)^2 - v_{\text{orbital}}(B)^2\right]
\end{align}

\begin{equation}\label{eq:zeeman}
  \boxed{\Delta E_{\text{Zeeman}} = \tfrac{1}{2}m_e\,v_{\text{magnetic}}^2}
\end{equation}

Since $v_{\text{magnetic}} \propto B$, this derives the Zeeman energy shift as a direct, mechanical consequence of budget reallocation.  It also predicts a fundamentally \textbf{non-linear Zeeman effect at extreme field strengths}---a falsifiable prediction of SDT.


%----------------------------------------------------------------------
\section{Conclusion}
%----------------------------------------------------------------------

The Unified Movement Budget is a universal conservation law governing particle dynamics.  From the single axiom that an electron vortex possesses a finite, conserved total velocity ($v_{\text{budget}} \approx 0.548\,c$), we have derived the mechanical origins of both relativistic time dilation and the quantum Zeeman effect.

These are not separate phenomena governed by disparate laws.  They are two faces of the same coin: the necessary reallocation of a particle's finite movement budget in response to external conditions.  This principle unifies key aspects of Special Relativity and Quantum Electromagnetism under a single, predictive geometric law.
